\section{Boundaries}

RCOUNT canvass arithmetic has four main boundaries.

First, it does not decide ballot eligibility. A package can record accepted,
counted, rejected, and cured ballot categories, but the legal decision to accept
or reject a ballot is made under jurisdiction-specific law.

Second, it does not replace canvassing-board minutes or certification records.
Status events should point to public authority and source references, but the
software verifies the presence and consistency of those records, not their legal
sufficiency.

Third, it does not require one national canvass workflow. The synthetic fixtures
are deliberately small and jurisdiction-neutral. Real adapters must preserve
state and local differences in cure periods, provisional rules, recount
thresholds, and reporting-unit conventions.

Fourth, it does not expose voter choices. Even when a future adapter ingests CVR
or ballot-manifest evidence, public RCOUNT packages must preserve the privacy
boundary described in V.0. Public event metadata also inherits that boundary:
small reporting units, rare write-ins, ballot-style combinations, timestamps,
or overly precise source references can create reidentification risk. Hashing a
source file provides tamper evidence; it does not make private source data safe
to publish.

Finally, V.1 is not an end-to-end cryptographic voting protocol. It supplies
replayable public arithmetic and evidence references. Later papers may add
privacy-safe inclusion proofs and CVR reconciliation, but the rule remains:
prove inclusion, never prove candidate choices.
